\section{Experimental Details}
\label{app:experimental-details}

% 为确保不同方法之间的比较具有可比性，所有方法均采用相同的 Qwen3-4B 基础模型、LoRA 配置和固定数据划分。
% 具体而言，LoRA 应用于所有线性层，秩为 $r\!=\!16$、缩放系数为 $\alpha\!=\!32$、dropout 为 $0.05$，
% 且设置 \texttt{bias=none}；数据划分采用 \texttt{split\_seed}=20260720，验证集比例为 $0.1$。
% 在训练阶段，所有方法均使用 $1\!\times\!10^{-4}$ 的学习率并采用余弦衰减，
% 同时设置 $0.03$ 的预热比例、3 个 epoch、$0.3$ 的最大梯度范数和 8 步梯度累积。
% SO、RS 和 ReaSon 采用标准 SFT，单设备训练和评估批大小均为 2；
% Mix、PaIR 和 RePaIR 采用配对训练目标，单设备训练批大小为 1，
% 每个训练样本包含配对的 DS 和 RF 视图，而评估批大小仍为 2。
% 此外，模型以 BF16 精度训练，采用 SDPA 注意力和梯度检查点（\texttt{use\_reentrant=False}）；
% 每个 epoch 结束时进行评估并保存 checkpoint，仅保留一个 checkpoint，并选择最终 checkpoint。
% 在推理阶段，采用 bfloat16 vLLM 和确定性的贪心解码；
% 其他训练与推理配置汇总于表~\ref{tab:train-infer-configs}。

To ensure comparability across methods, all methods use the same Qwen3-4B base model, 
LoRA configuration, and fixed data split. Specifically, LoRA is applied to all linear 
layers with rank $r\!=\!16$, scaling factor $\alpha\!=\!32$, a dropout rate of $0.05$, 
and \texttt{bias=none}; the data split is fixed using \texttt{split\_seed}=20260720 with 
a validation ratio of $0.1$. 
During training, we use a learning rate of $1\!\times\!10^{-4}$ 
with cosine decay, together with a warm-up ratio of $0.03$, three epochs, a maximum gradient 
norm of $0.3$, and eight gradient-accumulation steps. 
SO, RS, and ReaSon use standard SFT with per-device training and evaluation
batch sizes of 2. Mix, PaIR, and RePaIR instead use paired training objectives
with a per-device training batch size of 1, where each training example contains
paired DS and RF views, while the evaluation batch size remains 2.
Models are trained in BF16 using SDPA attention and gradient checkpointing 
(\texttt{use\_reentrant=False}); evaluation and checkpoint saving are performed after 
each epoch, only one checkpoint is retained, and the final checkpoint is selected. 
At inference time, we use bfloat16 vLLM with deterministic greedy decoding.

Other training and inference configurations for all methods are summarized in 
Table~\ref{tab:train-infer-configs}.

\newcolumntype{W}[1]{>{\hsize=#1\hsize\raggedright\arraybackslash}X}
\newcolumntype{N}[1]{>{\hsize=#1\hsize\centering\arraybackslash}X}
